\documentclass{article}
\usepackage{fullpage,upquote,times}
\sloppy

\begin{document}

\begin{verbatim}

Ideas for the new object-oriented constraint programming language

GOALS

integrate object-oriented programming with constraints
  address integration problems: support normal methods and messages;
  constraints on method results

predictable efficiency: efficient for basic built-in solvers

extensibility (add new solvers; cooperating solvers)



ARRAYS, SETS, AND LOOPS

Here's a version of Array sum that uses recursion to construct temp variables:

class Array
  def sum
    if self.empty?
      0
    else
      always { ans == self[0] + self[1..-1].sum }
      ans
    end
  end
end

or using a constraint version of inject:
class Array
  def sum
    constraint_inject(0) { |partial_sum,x| partial_sum+x }
  end
end

The second version assumes there is a constraint_inject method defined
on Enumerable ... this could be similar to the recursive Array sum method.

Both these should (in theory) interact correctly with other linear
constraints on the sum or on values.  For example, you could determine
an element given the sum and values for all the other elements:

a = Array.new(3)
a[0] = 0.0
a[1] = 0.0
a[2] = 0.0
always { a[0] == 10 }
always { a[2] == 20 }
always { a.sum == 50 }

This should be solvable and result in a[1] = 20

In conventional Ruby you'd just use a loop:

class Array
  def sum
    sum = 0
    each { |x| sum=sum+x }
    sum
  end
end

or inject.  However, I don't think that will work with constraints,
unless there is some way to automatically re-trigger the method.  Also
how would it work with a known sum and finding an element?

We also want constraints on the length of an array, and the whole array.
For arrays a and b:

always { a.length == 10 }
always { a == b}

If we use the Cooperating Solvers idea of multiple regions connected
by read-only variables, and require that constraints on an array
length be in a separate and upstream region from constraints on the
elements, and also constraints on the whole array (a==b) then we
should be able to have all these kinds.

Here is an example of constraints that wouldn't be allowed (because of
the cycle among the regions):

array a
a.length == a[0]
a.length <= 3
a.sum == 10
a[1] ==2

This has the solution a=[3,2,5], but finding that is probably beyond
the capabilities of the solvers we want.


WISH LIST FOR CONSTRAINTS ON DATA STRUCTURES

This wish list is very likely going to be beyond the capabilities of
the current solvers!

SETS
constraints involving 
  == subset subseteq member
and expressions involving union intersection - size

ARRAYS
  map member + (concat) == length [n] [n..m]

DICTIONARIES
  keys, values, [key] size etc

BINARY TREES
  depth, mirror, flatten, sum, etc

If we were going to provide these, we'd want constraints on finite
domains or specialized solvers, and perhaps backtracking.  (For
example, sets and arrays could be represented as CLP(R) lists, and
backtracking used to search for solutions.)


CHANGE OVER TIME

Midpoint line -- have an edit constraint and a loop

            always { (end1 + end2) == (midpoint*2) }

Need incrementality.

What is the relation to FRP etc?

Let's start with a simple version, just driven from a loop.  Later we
need to figure out how this works with events or some other mechanism.

m = MidpointLine.new(Point(0,0),Point(100,100))
(1..100).each { |i| m.end2.x = i }

How to express edit constraints??
  



THINGS TO ADD

read-only variables (should we use ?, or is that going to be a mess?)


THINGS TO FIX

versions of sum for arrays should work, also resistor with the current
 
class Battery
  # this seems unintuitive, and perhaps a common pattern ....
  def initialize(supply_voltage)
    super()
    @supply_voltage = supply_voltage
    always { lead1.voltage - lead2.voltage == @supply_voltage.stay(:required) }
  end

QUESTIONS

Syntax for priorities, both on stays and other constraints?

How do solvers communicate (use Cooperating Solvers design!!)
See old questions about priorities, incrementality, etc. in the VPRI memo.

What about Kaleidoscope concerns: identity constraints, constraint
constructors, constraints on class?

What about the things that motivated multi-methods (e.g. a constraint
a+b=c, with b and c known and a unknown)?  Or coercions?

The idea of constraints on messages, or more accurately on the results
of sending an arbitrary message, seems over general.  Should there be
restrictions?  Virtual part?  An alternative is to allow constraints
on arbitrary expressions (should evaluate to a boolean??) but to start
coming up with the restrictions that the programmer should follow,
even if they aren't checked.  

If always { } returns a ConstraintObject, what is the protocol for
ConstraintObject?  solver, variables, ... also error?  methods for
local propagation constraints?  linear expression for linear
constraints?


NAMES

The language needs a name!  Some ideas:

Kaleidoscope 2013
Mandala
Spargel
Brandenberg
Babelsberg
Tiergarten

\end{verbatim}
\end{document}
